\subject{Physics}
\subjectcode{5054}
\paper{2}
\variant{1}
\session{M/J}
\year{2025}

\begin{question}{1}
\markscheme
\row{1(a)}{Curve upwards (at start) with decreasing gradient labelled A}{B1}
\row{}{Horizontal line labelled B}{B1}
\row{1(b)}{Gradient/slope decreases -- increase in speed for same time is less at later times}{B1}
\row{1(c)(i)}{Air resistance / drag}{B1}
\row{1(c)(ii)}{Air resistance / upwards force increases (with increased speed)}{B1}
\row{}{Air resistance / upwards force equals / balances weight}{B1}
\row{1(d)}{Vector diagram showing 400 N and 100 N force at right angles with correct resultant and all directions shown}{B1}
\row{}{410 (N)}{B1}
\row{}{12--16 degrees}{B1}
\endms

\exemplar
(a) Speed-time graph: draw a curve that starts at the origin, rises steeply at first and then bends over with a steadily decreasing gradient, finally levelling off into a horizontal straight line. Mark A on the rising, curved part (where the speed is still increasing) and B on the flat horizontal part (constant / terminal speed).

(b) The gradient of a speed-time graph equals the acceleration. As time (and speed) increase, the curve becomes less steep, so its gradient decreases; therefore the acceleration decreases as the speed increases.

(c)(i) Air resistance (drag).

(c)(ii) As the skydiver speeds up, the air resistance acting upwards increases. Eventually the air resistance grows until it equals (balances) her weight. With the up and down forces equal, there is no resultant vertical force.

(d) Draw the 400 N force vertically downwards and the 100 N force horizontally to the right, joined tip-to-tail at right angles, each with an arrow; the resultant is the arrow from the start of the first to the end of the second.
$$ R = \sqrt{400^{2} + 100^{2}} = \sqrt{170000} = 410\ \text{N} $$
$$ \tan\theta = \frac{100}{400} \Rightarrow \theta = 14^{\circ}\ \text{to the vertical} $$
\endexemplar
\end{question}

\begin{question}{2}
\markscheme
\row{2(a)}{(Total) clockwise moments = anticlockwise moments}{M1}
\row{}{In equilibrium}{A1}
\row{2(b)(i)}{19.6 or 20 N}{B1}
\row{2(b)(ii)}{80 x 3 x (b)(i) = 90 x weight}{B1}
\row{}{52 (N) and weight of suitcase does not exceed maximum}{B1}
\row{2(c)(i)}{(Place) where weight of object (can be taken to) acts}{B1}
\row{2(c)(ii)}{Centre of gravity falls outside base / to the left of P}{B1}
\row{}{Weight causes anticlockwise moment / turning effect}{B1}
\endms

\exemplar
(a) For a body in equilibrium, the sum of the clockwise moments about any point equals the sum of the anticlockwise moments about that same point.

(b)(i) $$ W = mg = 2.0 \times 9.8 = 19.6\ \text{N} $$

(b)(ii) Taking moments about the pivot, the three bags of sugar at 80 cm balance the suitcase at 90 cm:
$$ (80)(3)(19.6) = (90)(W) $$
$$ W = \frac{80 \times 3 \times 19.6}{90} = 52\ \text{N} $$
Since 52 N is less than 67 N, the weight of the suitcase does not exceed the maximum allowed.

(c)(i) The point at which the whole weight of the object can be taken to act.

(c)(ii) When the slope is steep enough, the centre of gravity moves outside the base of the bus, to the left of point P. The weight of the bus then acts outside P, producing an anticlockwise moment (turning effect) about P, so the bus rotates and topples over.
\endexemplar
\end{question}

\begin{question}{3}
\markscheme
\row{3(a)}{Any three from: reduction in chemical energy store in battery; electrical work done by motor; mechanical work done on weight; increase in gravitational energy store of weight; increase in thermal energy (in motor or air); electrical heating (by current)}{B3}
\row{3(b)(i)}{Useful or output energy less than input energy}{B1}
\row{3(b)(ii)}{Total energy is constant}{B1}
\row{}{Energy from power supply = energy gained by weight + thermal energy (to air / due to friction)}{B1}
\row{3(c)}{(E =) power x time in any form}{C1}
\row{}{180 J}{A1}
\endms

\exemplar
(a) The chemical energy store in the battery decreases; the motor does electrical work; mechanical work is done raising the load, so the gravitational potential energy store of the load increases; and some energy is transferred to thermal energy stores (in the motor and the surrounding air) through electrical heating and friction. (any three)

(b)(i) The useful output energy (the gravitational potential energy gained by the load) is less than the input energy supplied, because some of the input energy is wasted as thermal energy.

(b)(ii) The total energy stays constant -- energy is only transferred, never created or destroyed. Here the energy supplied by the power supply equals the energy gained by the weight plus the thermal energy transferred to the surroundings (to the air / by friction).

(c) $$ E = \text{efficiency} \times \text{input energy} = 0.60 \times (P \times t) = 0.60 \times (15 \times 20) = 180\ \text{J} $$
\endexemplar
\end{question}

\begin{question}{4}
\markscheme
\row{4(a)(i)}{Faster particles escape (from water on skin)}{B1}
\row{}{Leaving behind slower particles}{B1}
\row{4(a)(ii)}{Any one of: evaporation occurs on the surface; boiling is at one temperature; boiling involves the formation of bubbles (in the liquid); boiling occurs inside the liquid}{B1}
\row{4(b)(i)}{Cold air sinks}{B1}
\row{}{Cold air is dense(r than hot air)}{B1}
\row{4(b)(ii)}{(Tf = Ti -) E/mc in any form or 14.8 or 15 seen}{C1}
\row{}{5.2 degrees C}{A1}
\endms

\exemplar
(a)(i) The faster-moving particles have enough energy to escape from the surface of the water. This leaves behind the slower-moving particles, so the average kinetic energy of the remaining liquid falls, which means its temperature falls and it cools (cooling the skin).

(a)(ii) Evaporation happens only at the surface of the liquid (and at any temperature), whereas boiling happens throughout the liquid at one fixed temperature, with bubbles forming inside it. (any one)

(b)(i) Air next to the cold coils is cooled and becomes denser than the surrounding warmer air, so it sinks. Warmer, less dense air rises to take its place and is cooled in turn, setting up a continuous convection current.

(b)(ii) $$ \Delta\theta = \frac{E}{mc} = \frac{160000}{3.6 \times 3000} = 14.8^{\circ}\text{C} $$
$$ T_f = 20.0 - 14.8 = 5.2^{\circ}\text{C} $$
\endexemplar
\end{question}

\begin{question}{5}
\markscheme
\row{5(a)(i)}{C and R marked correctly}{B1}
\row{5(a)(ii)}{Compressions are where pressure is high / particles close together / density is high}{B1}
\row{5(b)(i)}{0.75 (Hz)}{B1}
\row{5(b)(ii)}{Speed: no change}{B1}
\row{}{Wavelength: decreases}{B1}
\row{5(b)(iii)}{At least 3 wavefronts in the shallow region parallel to each other and on correct side of normal}{M1}
\row{}{Refracted towards the surface on correct side of surface and join wavefronts shown}{A1}
\row{5(c)}{Spreading out / bending at an edge / obstacle, or spreading out after passing through a (small) gap / hole}{B1}
\endms

\exemplar
(a)(i) Mark C at the centre of a region where the dots (particles) are bunched closest together (a compression), and mark R at the centre of a region where the dots are most spread out (a rarefaction).

(a)(ii) A compression is a region where the particles are close together and the pressure (and density) is high; a rarefaction is a region where the particles are spread apart and the pressure (and density) is low.

(b)(i) $$ f = \frac{\text{number of oscillations}}{\text{time}} = \frac{45}{60} = 0.75\ \text{Hz} $$

(b)(ii) Speed: no change (the speed depends on the water/depth, not the frequency). Wavelength: decreases -- since $v = f\lambda$ and $v$ is constant, a higher frequency gives a shorter wavelength.

(b)(iii) In the shallow region the wave travels more slowly, so the crests become closer together and change direction. Draw at least three crests in the shallow region, parallel to one another and closer-spaced than the incoming ones, refracted (turned) so that they join up correctly with the incoming crests on the right side of the boundary.

(c) Diffraction is the spreading out (and bending) of a wave as it passes the edge of an obstacle, or as it passes through a gap / opening.
\endexemplar
\end{question}

\begin{question}{6}
\markscheme
\row{6(a)(i)}{Voltmeter across X and ammeter in series anywhere in circuit}{B1}
\row{6(a)(ii)}{Any two from: X and R share voltage / p.d. of battery or e.m.f. of battery = p.d. across X + p.d. across R; as R varies, p.d. across X and R vary; if R increases then p.d. across R increases; if R increases then p.d. across X decreases}{B2}
\row{6(b)(i)}{Directly proportional}{B1}
\row{6(b)(ii)}{Resistance = 90 x 3.6 / 60 or 5.4 ohms seen}{B1}
\row{}{(I =) V/R in any form or with any resistance}{C1}
\row{}{1.7 A}{A1}
\row{6(b)(iii)}{As length increases current decreases}{C1}
\row{}{(Current) inversely proportional (to length)}{A1}
\row{6(c)}{Inversely proportional}{B1}
\endms

\exemplar
(a)(i) Connect the voltmeter in parallel across X (to measure the p.d. across X) and connect the ammeter in series anywhere in the circuit (to measure the current through X).

(a)(ii) X and R are in series, so the e.m.f. of the battery is shared between them: p.d. across X + p.d. across R = e.m.f. As the variable resistor R is changed, the share of the voltage changes -- increasing R increases the p.d. across R and so decreases the p.d. across X (and decreasing R does the reverse). (any two)

(b)(i) The resistance is directly proportional to the length of the wire.

(b)(ii) From the graph, 60 cm of wire has a resistance of 3.6 ohms. Since resistance is proportional to length:
$$ R = 3.6 \times \frac{90}{60} = 5.4\ \Omega $$
$$ I = \frac{V}{R} = \frac{9.0}{5.4} = 1.7\ \text{A} $$

(b)(iii) As the length of the wire increases, the current decreases. Because $R \propto \text{length}$ and $I = V/R$ at constant p.d., the current is inversely proportional to the length.

(c) The resistance of a wire is inversely proportional to its cross-sectional area.
\endexemplar
\end{question}

\begin{question}{7}
\markscheme
\row{7(a)}{Magnetic field produced by current (in primary)}{B1}
\row{}{(Iron) core transmits magnetic field to secondary coil, or magnetic field is changing / reversing / alternating}{B1}
\row{}{Induction of output voltage (in secondary coil)}{B1}
\row{7(b)(i)}{Off the scale (of the voltmeter) / too large (for voltmeter)}{B1}
\row{7(b)(ii)}{VS / VP = NP / NS in any form}{C1}
\row{}{96}{A1}
\row{7(c)(i)}{Sensible alternating trace either side of mid-line}{B1}
\row{7(c)(ii)}{Measure or mention of height of trace}{B1}
\row{}{Multiply height by Y-gain value}{B1}
\endms

\exemplar
(a) The alternating current in the primary coil produces a continually changing (alternating) magnetic field. The iron core carries and concentrates this magnetic field so that it passes through the secondary coil. The changing magnetic field through the secondary coil induces an alternating voltage across it (electromagnetic induction).

(b)(i) This is a step-up transformer with an output roughly twice the input, so a 6.4 V input would give about 12.8 V output. That is off the top of the voltmeter's 0--10 V scale (too large to read).

(b)(ii) The voltage ratio equals the turns ratio:
$$ \frac{V_S}{V_P} = \frac{N_S}{N_P} $$
From the data, $V_S/V_P = 2.4/1.2 = 2$, so:
$$ N_S = N_P \times \frac{V_S}{V_P} = 48 \times 2 = 96 $$

(c)(i) Draw a sensible alternating (sine-shaped) wave, centred on the mid-line (the 0 line) and symmetrical above and below it.

(c)(ii) Measure the height of the trace from the mid-line up to a peak, in divisions (squares) on the screen. Then multiply that number of divisions by the Y-gain setting (2.0 V per division) to obtain the maximum (peak) value of the output voltage.
\endexemplar
\end{question}

\begin{question}{8}
\markscheme
\row{8(a)}{Protons and neutrons}{C1}
\row{}{145 neutrons, 94 protons}{A1}
\row{8(b)(i)}{A neutron hits / absorbed by a (plutonium) nucleus}{B1}
\row{}{(At least) one released neutron causes another (Pu) nucleus to split / fission}{B1}
\row{8(b)(ii)}{Control rods absorb neutrons}{B1}
\row{}{Insert / lower rods (between fuel rods) to decrease reaction}{B1}
\row{8(c)}{Any two from: charge is different; mass is different; alpha particles made of 2 protons and/or 2 neutrons / beta particle is an electron; ionising effect; penetrating power}{B2}
\row{8(d)(i)}{Cloud chamber: closed container with source S inside or outside and some means of cooling, or spark counter: fine wire and plate with source nearby}{B1}
\row{}{Cloud chamber: presence of vapour / volatile liquid, or spark counter: (high) voltage between wire and plate or plates charged positive and negative}{B1}
\row{8(d)(ii)}{Ionisation (of air)}{B1}
\endms

\exemplar
(a) The nucleus contains protons and neutrons. For plutonium-239 there are 94 protons and $239 - 94 = 145$ neutrons.

(b)(i) A neutron is absorbed by a plutonium-239 nucleus, causing it to split (fission) and release two or three more neutrons. At least one of these released neutrons is absorbed by another plutonium nucleus, causing that one to fission too. The process repeats and builds up, giving a chain reaction.

(b)(ii) Control rods absorb neutrons. Lowering (inserting) the rods further between the fuel rods absorbs more neutrons, so fewer are left to cause fission and the reaction slows down. Raising (withdrawing) the rods absorbs fewer neutrons, so more are available to cause fission and the reaction speeds up.

(c) Difference 1: their charge is different -- an alpha particle has charge +2, whereas a beta particle has charge -1. Difference 2: their mass is different -- an alpha particle is relatively heavy (2 protons and 2 neutrons), while a beta particle is an electron of very small mass. (Differences in ionising effect or penetrating power are also acceptable.)

(d)(i) Either a cloud chamber or a spark counter is acceptable. Cloud chamber: draw a sealed, closed container with the radioactive source labelled S (inside, or just outside), containing a supersaturated vapour / volatile liquid, with a means of cooling (for example a cold base such as dry ice). Spark counter: draw a fine wire held close to a metal plate, with the source labelled S placed nearby, and a high voltage applied between the wire and the plate. Label the source S in whichever diagram is drawn.

(d)(ii) Ionisation of the air by the alpha particles -- the ions produced form the visible tracks in a cloud chamber, or trigger the sparks in a spark counter.
\endexemplar
\end{question}

\begin{question}{9}
\markscheme
\row{9(a)}{Protostar, stable star in that order for first two boxes, or red super giant, supernova, black hole in that order for last three boxes}{C1}
\row{}{Protostar, stable star, red super giant, supernova, black hole in that order}{A1}
\row{9(b)(i)}{Gravity / gravitational attraction mentioned}{B1}
\row{9(b)(ii)}{High temperature, or pressure of light / radiation}{B1}
\row{9(c)(i)}{Distance light travels in one year}{B1}
\row{9(c)(ii)}{8200 years}{B1}
\row{9(c)(iii)}{SN185 is not moving / little speed (away from Earth so no redshift)}{B1}
\row{}{SN2014J (is much further from Earth so) is moving away faster (so has substantial redshift), or larger redshift due to expansion of universe (with time)}{B1}
\row{9(d)}{By fusion}{B1}
\row{}{Supernova / explosion of star, or fusion mentioned in a star}{B1}
\endms

\exemplar
(a) cloud of dust and gas $\rightarrow$ protostar $\rightarrow$ stable star $\rightarrow$ red supergiant $\rightarrow$ supernova $\rightarrow$ black hole.

(b)(i) Gravity (gravitational attraction).

(b)(ii) The fusion reactions keep the core at a very high temperature, and this produces a high outward pressure (the pressure of the hot gas and of the radiation / light), which pushes outwards against gravity.

(c)(i) A light-year is the distance that light travels in one year.

(c)(ii) Light takes 8200 years to cross a distance of 8200 light-years.

(c)(iii) SN185 lies within our own Milky Way and is essentially not moving away from us, so its light shows no redshift. SN2014J is far outside our galaxy; because the Universe is expanding, more distant objects recede faster, so SN2014J is moving away rapidly and its light shows a large redshift.

(d) Heavier elements are built up by nuclear fusion inside stars, and the heaviest elements are formed and scattered into space when a massive star explodes as a supernova.
\endexemplar
\end{question}